% GENERATED FILE. Edit canonical lessons, then run npm run generate:books.
% canonical-source-hash: fnv1a64:edf3213576f4027b
% canonical-lessons: MR-C153-karda, MR-C153-soneri, MR-C153-cadhne, MR-C153-utarne, MR-C153-pohocne

\chapter{Grey, Golden, To climb, To get down, To arrive}
\label{ch:mr-grey-golden-to-climb-to-get-down-to-arrive}
\hlchaptermodality{153}

\begin{chapteropening}
\textbf{By the end of this chapter:} \emph{I can say grey, golden, to climb, to get down and to arrive.}

Complete the last lesson of chapter 153: I can say grey, golden, to climb, to get down and to arrive.
\end{chapteropening}

\section[\texorpdfstring{\mr{करडा} (karḍā)}{करडा (karḍā)}]{\mr{करडा} (karḍā) — grey}

\label{lesson:MR-C153-karda}



\begin{quote}
Before the new one: say the Marathi for purple, then the Marathi for orange.
\end{quote}

\subsection*{You'll want to know: \mr{करडा}}
\textbf{\mr{करडा}} — \emph{karḍā} — \textquotedblleft{}grey\textquotedblright{}.

\begin{cousinweb}
The origin of \textbf{\mr{करडा}} (\emph{karḍā}) is not settled. Neither CDIAL nor DEDR, searched by form and by meaning, has an entry for it.
\end{cousinweb}

\begin{grammarlens}[title={What you've built}]
One more word for this chapter.
\end{grammarlens}

\subsection*{Your turn}
\begin{itemize}
\raggedright
  \item \emph{Say it:} \emph{karḍā}
  \item \emph{Say it:} \emph{karḍā}, once more
  \item \emph{Say it:} say \emph{karḍā}
  \item \emph{Recall:} say the Marathi for purple, then the Marathi for orange, then say \emph{karḍā} again
\end{itemize}

\begin{tcolorbox}[breakable,colback=teal!4,colframe=teal!35!black,title={Before you move on}]
What does \mr{करडा} mean? (\textquotedblleft{}Grey\textquotedblright{}.) Say it once more.
\end{tcolorbox}
\section[\texorpdfstring{\mr{सोनेरी} (sonerī)}{सोनेरी (sonerī)}]{\mr{सोनेरी} (sonerī) — golden}

\label{lesson:MR-C153-soneri}



\begin{quote}
Before the new one: say the Marathi for orange, then the Marathi for grey.
\end{quote}

\subsection*{You'll want to know: \mr{सोनेरी}}
\textbf{\mr{सोनेरी}} — \emph{sonerī} — \textquotedblleft{}golden\textquotedblright{}.

\begin{cousinweb}
\textbf{\mr{सोनेरी}} (\emph{sonerī}) is a tadbhava, a word inherited and reshaped by sound change. It comes from \emph{*suvarṇadhara}, \textquotedblleft{}containing gold\textquotedblright{}, a form CDIAL reconstructs from its descendants (CDIAL 13522). It is built on \emph{suvarṇa}, \textquotedblleft{}gold\textquotedblright{}.
\end{cousinweb}

\begin{grammarlens}[title={What you've built}]
One more word for this chapter.
\end{grammarlens}

\subsection*{Your turn}
\begin{itemize}
\raggedright
  \item \emph{Say it:} \emph{sonerī}
  \item \emph{Say it:} \emph{sonerī}, once more
  \item \emph{Say it:} say \emph{sonerī}
  \item \emph{Recall:} say the Marathi for orange, then the Marathi for grey, then say \emph{sonerī} again
\end{itemize}

\begin{tcolorbox}[breakable,colback=teal!4,colframe=teal!35!black,title={Before you move on}]
What does \mr{सोनेरी} mean? (\textquotedblleft{}Golden\textquotedblright{}.) Say it once more.
\end{tcolorbox}
\section[\texorpdfstring{\mr{चढणे} (caḍhṇe)}{चढणे (caḍhṇe)}]{\mr{चढणे} (caḍhṇe) — to climb}

\label{lesson:MR-C153-cadhne}



\begin{quote}
Before the new one: say the Marathi for grey, then the Marathi for golden.
\end{quote}

\subsection*{You'll want to know: \mr{चढणे}}
\textbf{\mr{चढणे}} — \emph{caḍhṇe} — \textquotedblleft{}to climb\textquotedblright{}.

\begin{cousinweb}
\textbf{\mr{चढणे}} (\emph{caḍhṇe}) is a tadbhava, a word inherited and reshaped by sound change. It comes from \emph{*caḍhati}, \textquotedblleft{}rises\textquotedblright{}, a form CDIAL reconstructs from its descendants (CDIAL 4578).
\end{cousinweb}

\begin{grammarlens}[title={What you've built}]
One more word for this chapter.
\end{grammarlens}

\subsection*{Your turn}
\begin{itemize}
\raggedright
  \item \emph{Say it:} \emph{caḍhṇe}
  \item \emph{Say it:} \emph{caḍhṇe}, once more
  \item \emph{Say it:} say \emph{caḍhṇe}
  \item \emph{Recall:} say the Marathi for grey, then the Marathi for golden, then say \emph{caḍhṇe} again
\end{itemize}

\begin{tcolorbox}[breakable,colback=teal!4,colframe=teal!35!black,title={Before you move on}]
What does \mr{चढणे} mean? (\textquotedblleft{}To climb\textquotedblright{}.) Say it once more.
\end{tcolorbox}
\section[\texorpdfstring{\mr{उतरणे} (utarṇe)}{उतरणे (utarṇe)}]{\mr{उतरणे} (utarṇe) — to get down, to get off}

\label{lesson:MR-C153-utarne}



\begin{quote}
Before the new one: say the Marathi for golden, then the Marathi for to climb.
\end{quote}

\subsection*{You'll want to know: \mr{उतरणे}}
\textbf{\mr{उतरणे}} — \emph{utarṇe} — \textquotedblleft{}to get down, to get off\textquotedblright{}.

\begin{cousinweb}
\textbf{\mr{उतरणे}} (\emph{utarṇe}) is a tadbhava, a word inherited from Sanskrit and reshaped by sound change. It comes from \emph{uttarati}, \textquotedblleft{}passes over; alights, descends\textquotedblright{} (CDIAL 1770). CDIAL traces the sense from \textquotedblleft{}crosses over\textquotedblright{} to \textquotedblleft{}lands\textquotedblright{} to \textquotedblleft{}alights\textquotedblright{} and so \textquotedblleft{}descends\textquotedblright{}.
\end{cousinweb}

\begin{grammarlens}[title={What you've built}]
One more word for this chapter.
\end{grammarlens}

\subsection*{Your turn}
\begin{itemize}
\raggedright
  \item \emph{Say it:} \emph{utarṇe}
  \item \emph{Say it:} \emph{utarṇe}, once more
  \item \emph{Say it:} say \emph{utarṇe}
  \item \emph{Recall:} say the Marathi for golden, then the Marathi for to climb, then say \emph{utarṇe} again
\end{itemize}

\begin{tcolorbox}[breakable,colback=teal!4,colframe=teal!35!black,title={Before you move on}]
What does \mr{उतरणे} mean? (\textquotedblleft{}To get down, to get off\textquotedblright{}.) Say it once more.
\end{tcolorbox}
\section[\texorpdfstring{\mr{पोहोचणे} (pohocṇe)}{पोहोचणे (pohocṇe)}]{\mr{पोहोचणे} (pohocṇe) — to arrive}

\label{lesson:MR-C153-pohocne}



\begin{quote}
Before the new one: say the Marathi for to climb, then the Marathi for to get down.
\end{quote}

\subsection*{You'll want to know: \mr{पोहोचणे}}
\textbf{\mr{पोहोचणे}} — \emph{pohocṇe} — \textquotedblleft{}to arrive\textquotedblright{}.

\begin{cousinweb}
\textbf{\mr{पोहोचणे}} (\emph{pohocṇe}) is a tadbhava, a word inherited from Sanskrit and reshaped by sound change. It comes from \emph{prabhūta}, \textquotedblleft{}much, great\textquotedblright{} (CDIAL 8716). CDIAL glosses the Marathi verb \textquotedblleft{}to arrive\textquotedblright{}.
\end{cousinweb}

\begin{grammarlens}[title={What you've built}]
One more word for this chapter.
\end{grammarlens}

\subsection*{Your turn}
\begin{itemize}
\raggedright
  \item \emph{Say it:} \emph{pohocṇe}
  \item \emph{Say it:} \emph{pohocṇe}, once more
  \item \emph{Say it:} say \emph{pohocṇe}
  \item \emph{Recall:} say the Marathi for to climb, then the Marathi for to get down, then say \emph{pohocṇe} again
\end{itemize}

\begin{tcolorbox}[breakable,colback=teal!4,colframe=teal!35!black,title={Before you move on}]
What does \mr{पोहोचणे} mean? (\textquotedblleft{}To arrive\textquotedblright{}.) Say it once more.
\end{tcolorbox}
